\chapter{ساختار پایه و شبیه‌سازی وضعیت تلاطم}

در این فصل به بررسی ساختار پایهٔ شبیه‌ساز و نحوهٔ مدل‌سازی رفتار پردازنده برای ایجاد ترافیک حافظه می‌پردازیم. پیش از اضافه کردن هرگونه راهکار بهبود عملکرد، نیاز است تا معماری پایه به درستی مستقر شده و پدیدهٔ تلاطم\footnote{\lr{Thrashing}} که منجر به افت شدید عملکرد سامانه می‌شود به صورت دقیق شبیه‌سازی و اندازه‌گیری شود. برای این منظور از تولید الگوهای دسترسی متناوب به حافظه استفاده شده است.

\section{معماری پایه در شبیه‌ساز آکیتا}

شبیه‌ساز آکیتا\footnote{\lr{Akita}} یک چارچوب شبیه‌سازی رویدادمحور\footnote{\lr{Event-driven}} در معماری کامپیوتر است که به زبان گو\footnote{\lr{Go}} توسعه یافته است. این شبیه‌ساز از قطعات سخت‌افزاری مستقلی تشکیل شده است که هر کدام دارای پورت‌های ارتباطی مشخصی هستند و با ارسال پیام با یکدیگر ارتباط برقرار می‌کنند. آکیتا به طراحان اجازه می‌دهد تا زمان‌بندی دقیق چرخه‌های ساعت، تاخیر در صف‌های ارسال پیام و عملکرد قطعات حافظه را با دقت بالایی مدل‌سازی کنند. در این پروژه، از آکیتا برای ساخت یک سامانه پایه متشکل از یک پردازندهٔ تولیدکنندهٔ الگو، یک سطح از حافظهٔ نهان و یک کنترل‌کنندهٔ حافظهٔ اصلی استفاده شده است.\cite{akita}

\subsection{حافظهٔ نهان در شبیه‌ساز آکیتا}

پیکربندی حافظهٔ نهان سطح یک (\lr{L1}) در معماری پایه در فایل \lr{main.go} و با استفاده از سازندهٔ \lr{writeback} پیاده‌سازی شده است. مشخصات و ویژگی‌های عملکردی این حافظه به شرح زیر است:
\begin{itemize}
	\item \textbf{ظرفیت و اندازهٔ بلاک:} ظرفیت کل این حافظه برابر یک کیلوبایت (\lr{1 KB}) در نظر گرفته شده است. با توجه به تنظیم مقدار لگاریتم اندازهٔ بلاک برابر ۶ (\lr{Log2BlockSize = 6})، اندازهٔ هر بلاک داده در این حافظه ۶۴ بایت است. این یعنی حافظهٔ نهان در مجموع دارای ۱۶ بلاک فیزیکی است. ظرفیت و اندازهٔ حافظهٔ نهان در شکل \ref{cache} نمایش داده شده است.
	

\begin{figure}[ht]
	‫\centerimg{cache_memory}{14cm}
	\caption{ساختار و ظرفیت حافظهٔ نهان.}
	\label{cache}
\end{figure}

	\item \textbf{سیاست نگاشت:} درجهٔ انجمنی\footnote{\lr{Associativity}} این حافظه برابر ۱ تنظیم شده است. این مقدار بیانگر استفاده از روش نگاشت مستقیم\footnote{\lr{Direct-Mapped}} است. در این روش، هر آدرس حافظه دقیقا به یک خط مشخص در حافظهٔ نهان نگاشت می‌شود که زمینه را برای بروز نقصان‌های برخورد\footnote{\lr{Conflict Miss}} فراهم می‌کند. سیاست نگاشت به‌صورت تصویری در شکل \ref{mapping} آمده است.
	
\begin{figure}[ht]
	‫\centerimg{mapping}{12cm}
	\caption{سیاست نگاشت مستقیم.}
	\label{mapping}
\end{figure}
	
	\item \textbf{تاخیر دسترسی:} تاخیر اصابت\footnote{\lr{Hit Time}} برابر ۱ چرخهٔ ساعت است که به صورت مساوی (نیم چرخه) بین تاخیر دایرکتوری و تاخیر بانک‌های حافظه تقسیم شده است.
	\item \textbf{مدیریت درخواست‌های معلق:} این ماژول دارای ۴ مدخل در ثبات نگهداری وضعیت نقصان\footnote{\lr{Miss Status Holding Register - MSHR}} است. این ویژگی به حافظهٔ نهان اجازه می‌دهد در صورت وقوع نقصان، تا ۴ درخواست را به صورت هم‌زمان به سطوح پایین‌تر ارسال کرده و منتظر پاسخ بماند بدون این که متوقف شود.
\end{itemize}
علاوه بر موارد فوق، این ماژول از سیاست بازنویسی\footnote{\lr{Write-back}} استفاده می‌کند، به این معنا که تغییرات داده تا زمانی که خط مورد نظر از حافظه اخراج نشود، به حافظهٔ اصلی فرستاده نخواهد شد.

\subsection{حافظهٔ اصلی در شبیه‌ساز آکیتا}

برای مدل‌سازی حافظهٔ اصلی در حالت پایه، از یک کنترل‌کنندهٔ حافظهٔ ایده‌آل\footnote{\lr{Ideal Memory Controller}} استفاده شده است. این کنترل‌کننده پیچیدگی‌های داخلی ماژول‌های رم\footnote{\lr{RAM}} (مانند تازه کردن ردیف‌ها و تاخیرهای متغیر بانک‌ها) را نادیده گرفته و به جای آن یک مدل ساده و قابل پیش‌بینی ارائه می‌دهد. 
ظرفیت این حافظه ۶۴ کیلوبایت تعیین شده است که برای برنامه‌های محک\footnote{\lr{Benchmark}} این پروژه کاملا پاسخگو است. مهم‌ترین مشخصهٔ این بخش، تاخیر بسیار بالای آن است که برابر ۱۰۰ چرخهٔ ساعت (\lr{MEM\_LATENCY = 100}) تنظیم شده است. این تاخیر سنگین باعث می‌شود که هر بار داده‌ای در حافظهٔ نهان یافت نشود، پردازنده برای دریافت آن جریمهٔ زمانی بسیار بزرگی را متحمل شود. ساختار حافظهٔ اصلی در شکل \ref{main memory} نمایش داده شده است.

\begin{figure}[ht]
	‫\centerimg{main_memory}{8cm}
	\caption{ساختار و ظرفیت حافظهٔ اصلی.}
	\label{main memory}
\end{figure}

\subsection{نحوهٔ اتصال حافظهٔ نهان و حافظهٔ اصلی}

در شبیه‌ساز آکیتا، قطعات سخت‌افزاری از طریق سامانه اتصال مستقیم\footnote{\lr{Direct Connection}} به یکدیگر متصل می‌شوند. هر ماژول دارای پورت‌های بالا\footnote{\lr{Top Port}} (برای دریافت درخواست از قطعات نزدیک‌تر به پردازنده) و پایین\footnote{\lr{Bottom Port}} (برای ارسال درخواست به سطوح دورتر حافظه) است.
در معماری پایه، پورت حافظه (\lr{Mem}) در پردازنده به پورت بالای حافظهٔ نهان متصل می‌شود. سپس پورت پایین حافظهٔ نهان مستقیما به پورت بالای کنترل‌کنندهٔ حافظهٔ اصلی وصل می‌گردد. همچنین برای هدایت صحیح بسته‌های داده، از یک نگاشت‌گر پورت یگانه\footnote{\lr{Single Port Mapper}} استفاده شده است تا تمامی درخواست‌های خروجی از حافظهٔ نهان را بدون نیاز به مسیریابی پیچیده، مستقیما به سمت حافظهٔ اصلی هدایت کند. تمامی این اتصالات در فرکانس ۱ گیگاهرتز همگام‌سازی شده‌اند.

\subsection{تغییرات اعمال‌شده در کد شبیه‌ساز}
برای ارزیابی عملکرد حافظهٔ نهان، قابلیت جمع‌آوری آمار به ماژول \lr{Write-Back Cache} در شبیه‌ساز آکیتا اضافه شد. بدین منظور، در فایل \lr{mem/cache/writeback/writebackcache.go}، شمارنده‌هایی برای ثبت تعداد اصابت خواندن \footnote{\lr{Read Hit}}، 
نقصان خواندن\footnote{\lr{Read Miss}}، 
اصابت نوشتن\footnote{\lr{Write Hit}}، 
نقصان نوشتن\footnote{\lr{Write Miss}}، 
اخراجی‌ها\footnote{\lr{Evictions}} و 
بازنویسی‌ها به ساختار حافظهٔ نهان افزوده شد. سپس در فایل‌های مربوط به پردازش درخواست‌ها و جایگزینی بلوک‌ها، این شمارنده‌ها در محل مناسب به‌روزرسانی شدند تا آمار هر دسترسی به‌صورت دقیق ثبت شود. علاوه بر این، برای پیاده‌سازی حافظهٔ نهان قربانی نیز در فایل \lr{victimcache.go} شمارنده‌های مربوط به اصابت‌های حافظهٔ نهان قربانی\footnote{\lr{Victim Cache Hits}}،
 نقصان‌های حافظهٔ نهان قربانی\footnote{\lr{Victim Cache Misses}}
  و ترافیک حافظه\footnote{\lr{Memory Traffic}} 
اضافه شد تا عملکرد این سطح از حافظه نیز قابل ارزیابی باشد. در نهایت، در فایل \lr{tracecpu.go} واقع در پروژهٔ اصلی، این آمارها پس از اجرای برنامهٔ محک نمایش داده شده و با استفاده از آن‌ها نرخ \lr{Hit/Miss} و همچنین \lr{Average Memory Access Time (AMAT)} برای هر دو حالت، با و بدون حافظهٔ نهان قربانی، محاسبه و گزارش می‌شود.

\section{برنامهٔ محک و ایجاد تلاطم}

برای سنجش عملکرد حافظهٔ نهان و ارزیابی راه‌حل‌های پیشنهادی، نیاز به تولید ترافیک حافظه داریم. ماژول \lr{TraceCPU} در این پروژه وظیفه دارد به جای اجرای یک کد اسمبلی واقعی، مستقیما درخواست‌های خواندن و نوشتن را به آدرس‌های از پیش تعیین شده ارسال کند. در این بخش دو الگوی متفاوت دسترسی به حافظه به نام‌های الگوی کوچک (\lr{small trace}) و الگوی بزرگ (\lr{large trace}) طراحی شده است که هدف اصلی آن‌ها ایجاد ترافیک متناوب و شبیه‌سازی بدترین حالت برای حافظهٔ نهان نگاشت مستقیم، یعنی پدیدهٔ تلاطم است.

\subsection{آزمون اول}

آزمون اول با استفاده از الگوی \lr{smallTrace} اجرا می‌شود و هدف آن بررسی رفتار حافظه با تعداد کمی درخواست، اما با نرخ برخورد بالا است. این الگو به سه فاز اصلی تقسیم می‌شود:
\begin{enumerate}
	\item \textbf{فاز پر کردن حافظه:} در ابتدا ۱۶ درخواست نوشتن متوالی با فاصلهٔ ۶۴ بایت (از آدرس \lr{0x0000} تا \lr{0x03C0}) به حافظه ارسال می‌شود. چون ظرفیت حافظهٔ نهان دقیقا ۱۶ بلاک ۶۴ بایتی است، این فاز کل ظرفیت حافظهٔ نهان را پر می‌کند و به طور طبیعی ۱۶ نقصان اولیه\footnote{\lr{Cold Miss}} ایجاد می‌کند.
	\item \textbf{فاز تایید جایگذاری:} سپس پردازنده ۶ بلاک اول را دوباره می‌خواند که تمامی این خواندن‌ها منجر به اصابت موفقیت‌آمیز در حافظهٔ نهان می‌شوند.
	\item \textbf{فاز ایجاد تلاطم:} در این مرحله پردازنده درخواست نوشتن در آدرس \lr{0x0400} را ارسال می‌کند. فاصلهٔ آدرس \lr{0x0400} (معادل ۱۰۲۴ بایت) با آدرس \lr{0x0000} دقیقا برابر با ظرفیت کل حافظهٔ نهان (۱ کیلوبایت) است. به دلیل ساختار نگاشت مستقیم، هر دو آدرس به یک خط فیزیکی در حافظهٔ نهان ختم می‌شوند. با واکشی \lr{0x0400}، دادهٔ \lr{0x0000} اخراج می‌شود. بلافاصله در دستور بعدی، پردازنده دوباره آدرس \lr{0x0000} را درخواست می‌کند که باعث اخراج \lr{0x0400} می‌گردد. این فرایند که داده‌ها پیش از استفادهٔ مجدد از حافظه اخراج می‌شوند، پدیدهٔ تلاطم نام دارد. این الگو برای بلاک‌های بعدی نیز تکرار می‌شود تا ترافیک بیهودهٔ زیادی روانهٔ حافظهٔ اصلی گردد.
\end{enumerate}

\subsection{آزمون دوم}

آزمون دوم که با متغیر \lr{largeTrace} شناخته می‌شود، برای شبیه‌سازی یک بار کاری سنگین و طولانی‌تر طراحی شده است تا تاثیر تاخیر حافظه بر زمان کل اجرا را به شکل واضح‌تری نشان دهد.
در این آزمون، پردازنده پس از پر کردن اولیهٔ حافظه، وارد یک حلقهٔ تکرار ۳۰ تایی می‌شود. درون این حلقهٔ اصلی، حلقه‌های داخلی دیگری وجود دارند که ۱۰ بار به طور متناوب بین آدرس‌های کاملا متداخل درخواست ارسال می‌کنند (برای مثال تناوب مداوم بین آدرس‌های \lr{0x0000} و \lr{0x0400}، و سپس \lr{0x0440} و \lr{0x0040}). 
این حجم از درخواست‌های متناوب به آدرس‌هایی که با یکدیگر برخورد دارند، باعث می‌شود نرخ اصابت در حافظهٔ نهان سطح یک به شدت سقوط کند و سامانه در یک چرخهٔ باطل از اخراج و واکشی مجدد گرفتار شود. این آزمون بیش از ۶۰۰ دسترسی به حافظه را شبیه‌سازی می‌کند و در معماری پایه، تقریبا تمام این درخواست‌ها به عنوان نقصان برخورد شناخته شده و با تحمیل تاخیر ۱۰۰ چرخه‌ای، باعث کندی شدید زمان نهایی اجرای برنامه می‌شوند.